% conclusions.tex (국문 v14, 2026-10-20): 이 연구의 기여만 적고, 다른 곳에서 이미 확립된 진술은 고찰에 둔다.
\section{결론}\label{sec:conclusions}
이 연구는 활동층 두께를 지점-연도 설계로 평가하고 지도화하였다. 모든 라벨, 기후 입력, 예측, 점수는 명시된 한 해의 관측 지점 하나를 가리킨다. 세 가지 결과가 이 연구에 고유하다. 첫째, 지역 라벨로 재보정한 Stefan 모델에 그 잔차를 학습한 gradient boosting 을 더한 모델은 공간 블록 검증에서 Stefan 오차를 레나 델타에서 \ndrefLeblkRoabs{} cm, 캐나다에서 \ndrefCablkRoabs{} cm 줄였고, 이득은 조건부였다. 토양 · 지형 공변량이 필요하였고, 지역 Stefan 계수가 지점들을 잘못 나타내는 곳에서 가장 컸다. 둘째, 물리 기준값은 전이성을 보존하였다. 대상 지역 라벨 없이 물리 잔차 모델은 Stefan 모델 수준을 유지한 반면 직접 ML은 레나 델타에서 Stefan 오차를 \ndrefLeregDzabs{} cm 넘었고, 지역 라벨 5--25개로 레나 델타와 외부 CALM 지역의 물리 잔차 오차가 줄었다. 셋째, 모든 모델과 공개 제품 5종을 같은 셀 · 같은 해에서 채점한 결과 ML 제품의 정확도는 학습 지점의 위치를 따랐다. 레나 델타에서 제품은 얇은 활동층을 나타내지 못하였고(관측과의 상관 \nrLeWei{}, \nrLeLiu), 그 지역 라벨 없이 학습한 모델보다 부정확하였다. 이 워크플로는 세 지역의 2011--2024년 연도별 1 km 지도를 적용 가능 영역과 예측 구간과 함께 제공하며, 다른 영구동토 지역에 이 방법을 적용할 수 있는지, 라벨이 몇 개나 필요한지를 판단하는 데 필요한 평가 프로토콜과 라벨 수 근거를 함께 제공한다.
